\section{结束语}

本文报告的结果为非阿贝尔规范理论的本质以及格点 QCD 的连续极限提供了若干新的认识。
它们显然在多个方面仍不完整，并引出一些有趣的问题。特别是，要对其重正化性质获得结
构性的理解，很可能需要对威尔逊流作全阶的微扰分析。

第 2 节中的一圈计算提示，无论夸克味数多少，在正流时间得到的场都是重正化场。迄今
这个问题只在纯规范理论中被数值地研究过；但令人鼓舞的是可以指出：作为规范 Ward
恒等式的推论，该猜想对 QED 以及任意带电物质多重态都成立。

变换后的泛函积分~\eqref{eq:fint} 的一个引人入胜之处是它由光滑场所支配。由于场变换
~\eqref{eq:fieldtrafo} 是可逆的，这个积分仍然编码了理论所描述的全部物理。不过，
理论的某些性质（例如拓扑扇区的分割）在这种表述下变得特别透明，而其在高能下的行为
则用基本规范场来讨论更为合适。
